\documentclass[10pt,a4paper]{amsart}
\usepackage[margin=19mm]{geometry}
\usepackage{amsmath,amssymb,mathtools}
\usepackage[scr=rsfs,cal=euler]{mathalfa}
\usepackage{xfrac}
\usepackage[unicode,hidelinks,bookmarks=false]{hyperref}
\newcommand{\ie}{i.e.\ }
\newcommand{\cf}{cf.\ }
\newcommand{\resp}{respectively\ }
\newcommand{\Subsection}{\subsection}
\providecommand{\nobreakdash}{\nobreak\hbox{-}}
\setlength{\emergencystretch}{2em}
\multlinegap0pt
\usepackage{bm}
\usepackage{bbm}
\usepackage{dsfont}
\usepackage{xspace}
\usepackage{cancel}
\usepackage{mathtools}
\usepackage{amsthm}
\newtheorem{lemma}{Lemma}
\title[A Unified Framework for Rerandomization using Quadratic Form]{A Unified Framework for Rerandomization using Quadratic Forms (Lemma 1)}
\author{Kyle Schindl, Zach Branson}
\date{Source: arXiv:2403.12815v3 (CC BY 4.0)}
\begin{document}
\maketitle

\noindent\emph{Source and licence.} A Unified Framework for Rerandomization using Quadratic Forms. Version arXiv:2403.12815v3 (CC BY 4.0), distributed under the Creative Commons Attribution 4.0 International licence, as recorded in the arXiv licence metadata for this exact version and shown on the abstract page https://arxiv.org/abs/2403.12815v3. Full rights notice: licenses/arxiv-2403.12815-CC-BY-4.0.txt.\par\smallskip

\noindent\textit{Editorial scope.} Counted unit: the Lemma printed as Lemma 1 of the article (source0.tex lines 1404--1410, source label \texttt{Gamma\_cond\_exp}), delivered together with the article's own complete original proof of it (source0.tex lines 1411--1452, one \texttt{proof} environment of 42 lines). No mathematics is written, repaired or re-derived: statement and proof are the article's own text line for line. Required context retained uncounted: none. Every definition, display and label of those lines is retained unchanged. Omitted from this package and not redistributed: 1--1403, 1453--1653. Nothing inside a retained excerpt is omitted or altered: every retained body is delivered exactly as the article's lines are written, so no omission receipt of any kind accompanies this package. Citations: the retained excerpts carry no citation command at all, so no bibliography entry of the article is transcribed and no reference is redistributed. Reference adaptation: none was needed and none was made; every reference inside the delivered unit resolves inside it, so no unresolved reference or citation is produced. Printed slips kept literal: no printed word, symbol, hypothesis or number of the counted statement or of its proof was altered. Layout and class: the article's own class is \texttt{article} with packages including fancyhdr, hyperref, amsfonts, epsfig, graphicx, afterpage, pgfplots, bm, tikz, nicefrac, comment, bbm, mathtools, amsmath; the wrapper is the lane's pinned \texttt{amsart} editorial wrapper at 10pt on A4, which restates the theorem-like environments the retained text needs and adds the article's own macro definitions that the retained text needs (none), with extra wrapper packages (bm, bbm, dsfont, xspace, cancel, mathtools). The delivered numbering is the unit's own. Assets: no retained span contains a figure environment, a graphics-inclusion command, a PDF-page inclusion, a table or a TikZ picture; no image file is copied into this package, the package declares no asset dependency and no image file is redistributed with it. Licence: version arXiv:2403.12815v3 of this article is distributed under the Creative Commons Attribution 4.0 International licence, as recorded in the arXiv licence metadata for this exact version and shown on the abstract page https://arxiv.org/abs/2403.12815v3. A full rights notice is delivered at licenses/arxiv-2403.12815-CC-BY-4.0.txt. Characters: every retained excerpt is pure ASCII, so the delivered engine, XeLaTeX with its default Latin Modern text font, reports no missing character.
\begin{lemma} \label{Gamma_cond_exp}
    Suppose $X \sim \text{Gamma}(\alpha, \beta)$. Then, for some $a > 0$,
    \begin{align*}
        \mathbb{E}\left[X \mid X < a \right] = \alpha \beta \left(\frac{F(a; \alpha + 1, \beta)}{F(a; \alpha, \beta)}\right)
    \end{align*}
    where $F(a; \cdot, \cdot)$ is the cumulative distribution function of $X \sim \text{Gamma}(\cdot, \cdot)$.
\end{lemma}

\begin{proof}[\textbf{Proof:}] First, note that we parameterize $X \sim \text{Gamma}(\alpha, \beta)$ such that $X$ has the density
\begin{align*}
    f(x; \alpha, \beta) = \frac{x^{\alpha - 1} e^{- x / \beta}}{\beta^\alpha \Gamma(\alpha)}.
\end{align*}
Then, observe that
\begin{align*}
    \mathbb{E}\left[X \mid X < a \right] &= \frac{\mathbb{E}\left[X \cdot \mathbbm{1}(X < a) \right]}{\mathbb{P}(X < a)} = \frac{\int^a_0 x \left( \frac{x^{\alpha - 1} e^{- x / \beta}}{\beta^\alpha \Gamma(\alpha)}  \right) dx}{\int^a_0 \frac{x^{\alpha - 1} e^{- x / \beta}}{\beta^\alpha \Gamma(\alpha)} dx} = \frac{\frac{1}{\Gamma(\alpha)}\int^a_0  \frac{x^{\alpha} e^{- x / \beta}}{\beta^\alpha }  dx}{\frac{1}{\Gamma(\alpha)}\int^a_0 \frac{x^{\alpha - 1} e^{- x / \beta}}{\beta^\alpha } dx}.
\end{align*}
Next, we can write this expression as a ratio of lower incomplete gamma functions, $\gamma(s, x) = \int^x_0 t^{s-1} e^{-t} dt$, since
\begin{align*}
  \frac{\int^a_0  \frac{x^{\alpha} e^{- x / \beta}}{\beta^\alpha }  dx}{\int^a_0 \frac{x^{\alpha - 1} e^{- x / \beta}}{\beta^\alpha } dx}  &=    \frac{\int^a_0 (x / \beta)^{\alpha} e^{- x / \beta}  dx}{\frac{1}{\beta}\int^a_0 (x / \beta)^{\alpha - 1} e^{- x / \beta} dx} = \beta \frac{\int^{a / \beta}_0 y^{\alpha} e^{- y}  dy}{\int^{a / \beta}_0 y^{\alpha - 1} e^{- y} dy} = \beta \cdot \frac{\gamma(\alpha + 1, a / \beta)}{\gamma(\alpha, a / \beta)}
\end{align*}
where in the second to last equality we let $y = x / \beta$. We can further simplify this expression by noting that $\gamma(\alpha + 1, a / \beta) = \alpha \gamma( \alpha, a / \beta) - (a / \beta)^\alpha e^{- a / \beta}$. Then, we can see that
\begin{align*}
    \beta \cdot \frac{\gamma(\alpha + 1, a / \beta)}{\gamma(\alpha, a / \beta)} = \beta \left(\alpha - \frac{(a / \beta)^\alpha e^{- a / \beta}}{\gamma(\alpha, a / \beta)}\right) = \alpha \beta \left(1 - \frac{(a / \beta)^\alpha e^{- a / \beta}}{\alpha \gamma(\alpha, a / \beta)}\right).
\end{align*}
Finally, using the fact that $\Gamma(\alpha + 1) = \alpha \Gamma(\alpha)$, it follows that
\begin{align*}
     \frac{(a / \beta)^\alpha e^{- a / \beta}}{\alpha \gamma(\alpha, a / \beta)} = \frac{\frac{a^\alpha e^{- a / \beta}}{\beta^\alpha \Gamma(\alpha + 1)} }{\frac{\alpha \gamma(\alpha, a / \beta)}{\Gamma(\alpha + 1)}} = \frac{\beta \left(\frac{a^\alpha e^{- a / \beta}}{\beta^{\alpha+1} \Gamma(\alpha + 1)} \right) }{\frac{\gamma(\alpha, a / \beta)}{\Gamma(\alpha)}} = \beta \left(\frac{f(a; \alpha + 1, \beta)}{F(a; \alpha, \beta)} \right).
\end{align*}
Putting everything together, we have that
\begin{align*}
    \mathbb{E}\left[X \mid X < a \right] = \alpha \beta \left(1 - \beta \cdot \frac{f(a; \alpha + 1, \beta)}{F(a; \alpha, \beta)}\right).
\end{align*}
Note that this makes sense intuitively, since $\mathbb{E}[X] = \alpha \beta$, so we can see that as $a \to \infty$ we recover the unconditional mean. We can also see that the mean of $X$ is reduced by $\beta \cdot \frac{f(a; \alpha + 1, \beta)}{F(a; \alpha, \beta)}$. However, it is still possible to simplify this conditional expectation. To do so, observe that
\begin{align*}
    F(a; \alpha, \beta) - \beta f(a ; \alpha + 1, \beta) &= \frac{\gamma(\alpha, a / \beta)}{\Gamma(\alpha)} - \beta \left(\frac{a^\alpha e^{- a / \beta}}{\beta^{\alpha+1} \Gamma(\alpha + 1)} \right) \\
    &= \frac{\alpha \gamma(\alpha, a / \beta)}{\Gamma(\alpha + 1)} - \frac{\beta a^\alpha e^{- a / \beta}}{\beta^{\alpha+1} \Gamma(\alpha + 1)} \\
    &= \frac{\beta^{\alpha + 1} \alpha \gamma(\alpha, a / \beta) - \beta a^\alpha e^{- a / \beta}}{\beta^{\alpha + 1}\Gamma(\alpha + 1)} \\
    &= \frac{\beta^{\alpha + 1} \left( \gamma(\alpha + 1, a / \beta) + (a / \beta)^\alpha e^{- a / \beta}\right) - \beta a^\alpha e^{- a / \beta}}{\beta^{\alpha + 1}\Gamma(\alpha + 1)} \\
    &= \frac{ \gamma(\alpha + 1, a / \beta)}{\Gamma(\alpha + 1)} + \frac{\beta a^\alpha e^{- a / \beta} - \beta a^\alpha e^{- a / \beta}}{\beta^{\alpha + 1}\Gamma(\alpha + 1)} \\
    &= F(a; \alpha + 1, \beta)
\end{align*}
where the fourth equality follows since $\alpha \gamma( \alpha, a / \beta) = \gamma(\alpha + 1, a / \beta) + (a / \beta)^\alpha e^{- a / \beta}$. Thus,
\begin{align*}
    1 - \beta \cdot \frac{f(a; \alpha + 1, \beta)}{F(a; \alpha, \beta)} = \frac{F(a; \alpha, \beta) - \beta f(a; \alpha + 1, \beta)}{F(a; \alpha, \beta)} = \frac{F(a; \alpha + 1, \beta)}{F(a; \alpha, \beta)}
\end{align*}
which yields a final conditional expectation of
\begin{align*}
    \mathbb{E}\left[X \mid X < a \right] = \alpha \beta \left(\frac{F(a; \alpha + 1, \beta)}{F(a; \alpha, \beta)}\right)
\end{align*}
\end{proof}

\end{document}
